\documentclass[12pt]{article}
\usepackage{amsmath, amssymb, geometry}
\geometry{a4paper, margin=1in}

\begin{document}

\title{Homework 15.7: Triple Integrals in Cylindrical Coordinates\\ \Large Problem Sets and Solutions}
\author{}
\date{}
\maketitle

\section*{Problem Sets}

\textbf{Problem 1:} Plot the point whose cylindrical coordinates are given. Then find the rectangular coordinates of the point: $(4, \frac{\pi}{6}, 2)$.
\\
\textbf{Solution:} We are given $(r, \theta, z) = (4, \frac{\pi}{6}, 2)$.
\begin{align*}
x &= r \cos(\theta) = 4 \cos\left(\frac{\pi}{6}\right) = 4 \left(\frac{\sqrt{3}}{2}\right) = 2\sqrt{3} \\
y &= r \sin(\theta) = 4 \sin\left(\frac{\pi}{6}\right) = 4 \left(\frac{1}{2}\right) = 2
\end{align*}
The $z$-coordinate remains the same: $z = 2$. 
\\
\textbf{Final Answer:} $(2\sqrt{3}, 2, 2)$
\vspace{0.5cm}

\textbf{Problem 2:} Find the rectangular coordinates of the point given in cylindrical coordinates: $(5, \frac{3\pi}{4}, -1)$.
\\
\textbf{Solution:} We are given $(r, \theta, z) = (5, \frac{3\pi}{4}, -1)$.
\begin{align*}
x &= 5 \cos\left(\frac{3\pi}{4}\right) = 5 \left(-\frac{\sqrt{2}}{2}\right) = -\frac{5\sqrt{2}}{2} \\
y &= 5 \sin\left(\frac{3\pi}{4}\right) = 5 \left(\frac{\sqrt{2}}{2}\right) = \frac{5\sqrt{2}}{2}
\end{align*}
The $z$-coordinate is $-1$.
\\
\textbf{Final Answer:} $\left(-\frac{5\sqrt{2}}{2}, \frac{5\sqrt{2}}{2}, -1\right)$
\vspace{0.5cm}

\textbf{Problem 3:} Find the rectangular coordinates of the point given in cylindrical coordinates: $(2, \frac{4\pi}{3}, 4)$.
\\
\textbf{Solution:} We are given $(r, \theta, z) = (2, \frac{4\pi}{3}, 4)$.
\begin{align*}
x &= 2 \cos\left(\frac{4\pi}{3}\right) = 2 \left(-\frac{1}{2}\right) = -1 \\
y &= 2 \sin\left(\frac{4\pi}{3}\right) = 2 \left(-\frac{\sqrt{3}}{2}\right) = -\sqrt{3}
\end{align*}
The $z$-coordinate is $4$.
\\
\textbf{Final Answer:} $(-1, -\sqrt{3}, 4)$
\vspace{0.5cm}

\textbf{Problem 4:} Find the rectangular coordinates of the point given in cylindrical coordinates: $(8, \frac{\pi}{2}, -5)$.
\\
\textbf{Solution:} We are given $(r, \theta, z) = (8, \frac{\pi}{2}, -5)$.
\begin{align*}
x &= 8 \cos\left(\frac{\pi}{2}\right) = 8(0) = 0 \\
y &= 8 \sin\left(\frac{\pi}{2}\right) = 8(1) = 8
\end{align*}
The $z$-coordinate is $-5$.
\\
\textbf{Final Answer:} $(0, 8, -5)$
\vspace{0.5cm}

\textbf{Problem 5:} Change from rectangular to cylindrical coordinates: $(2\sqrt{3}, 2, 5)$.
\\
\textbf{Solution:} Given $(x, y, z) = (2\sqrt{3}, 2, 5)$.
\begin{align*}
r &= \sqrt{x^2 + y^2} = \sqrt{(2\sqrt{3})^2 + 2^2} = \sqrt{12 + 4} = \sqrt{16} = 4 \\
\tan(\theta) &= \frac{y}{x} = \frac{2}{2\sqrt{3}} = \frac{1}{\sqrt{3}}
\end{align*}
Since $x > 0$ and $y > 0$, the point is in the first quadrant. Thus, $\theta = \frac{\pi}{6}$.
\\
\textbf{Final Answer:} $(4, \frac{\pi}{6}, 5)$
\vspace{0.5cm}

\textbf{Problem 6:} Change from rectangular to cylindrical coordinates: $(-3, 3, 1)$.
\\
\textbf{Solution:} Given $(x, y, z) = (-3, 3, 1)$.
\begin{align*}
r &= \sqrt{(-3)^2 + 3^2} = \sqrt{9 + 9} = \sqrt{18} = 3\sqrt{2} \\
\tan(\theta) &= \frac{y}{x} = \frac{3}{-3} = -1
\end{align*}
Since $x < 0$ and $y > 0$, the point is in the second quadrant. The standard angle is $\theta = \frac{3\pi}{4}$.
\\
\textbf{Final Answer:} $(3\sqrt{2}, \frac{3\pi}{4}, 1)$
\vspace{0.5cm}

\textbf{Problem 7:} Change from rectangular to cylindrical coordinates: $(0, -5, -2)$.
\\
\textbf{Solution:} Given $(x, y, z) = (0, -5, -2)$.
\begin{align*}
r &= \sqrt{0^2 + (-5)^2} = 5 \\
\tan(\theta) &= \frac{-5}{0} \text{ (undefined)}
\end{align*}
The point lies on the negative $y$-axis, which corresponds to $\theta = \frac{3\pi}{2}$.
\\
\textbf{Final Answer:} $(5, \frac{3\pi}{2}, -2)$
\vspace{0.5cm}

\textbf{Problem 8:} Change from rectangular to cylindrical coordinates: $(-\sqrt{2}, -\sqrt{6}, 0)$.
\\
\textbf{Solution:} Given $(x, y, z) = (-\sqrt{2}, -\sqrt{6}, 0)$.
\begin{align*}
r &= \sqrt{(-\sqrt{2})^2 + (-\sqrt{6})^2} = \sqrt{2 + 6} = \sqrt{8} = 2\sqrt{2} \\
\tan(\theta) &= \frac{-\sqrt{6}}{-\sqrt{2}} = \sqrt{3}
\end{align*}
Since $x < 0$ and $y < 0$, the point is in the third quadrant. Thus, $\theta = \frac{4\pi}{3}$.
\\
\textbf{Final Answer:} $(2\sqrt{2}, \frac{4\pi}{3}, 0)$
\vspace{0.5cm}

\textbf{Problem 9:} Write the following equation in cylindrical coordinates: $x^2 + y^2 + 2z^2 = 16$.
\\
\textbf{Solution:} Substitute $x^2 + y^2 = r^2$:
\begin{align*}
r^2 + 2z^2 = 16
\end{align*}
\textbf{Final Answer:} $r^2 + 2z^2 = 16$
\vspace{0.5cm}

\textbf{Problem 10:} Write the following equation in cylindrical coordinates: $z = 3x^2 - 3y^2$.
\\
\textbf{Solution:} Substitute $x = r\cos(\theta)$ and $y = r\sin(\theta)$:
\begin{align*}
z &= 3(r\cos(\theta))^2 - 3(r\sin(\theta))^2 \\
z &= 3r^2(\cos^2(\theta) - \sin^2(\theta))
\end{align*}
Using the double angle identity $\cos(2\theta) = \cos^2(\theta) - \sin^2(\theta)$:
\\
\textbf{Final Answer:} $z = 3r^2\cos(2\theta)$
\vspace{0.5cm}

\textbf{Problem 11:} Write the following equation in cylindrical coordinates: $x + 2y + 3z = 6$.
\\
\textbf{Solution:} Substitute $x = r\cos(\theta)$ and $y = r\sin(\theta)$:
\begin{align*}
r\cos(\theta) + 2r\sin(\theta) + 3z = 6
\end{align*}
\textbf{Final Answer:} $r(\cos(\theta) + 2\sin(\theta)) + 3z = 6$
\vspace{0.5cm}

\textbf{Problem 12:} Identify the surface whose equation is given by $r^2 + z^2 = 36$.
\\
\textbf{Solution:} We know the conversion formula $r^2 = x^2 + y^2$. Substituting this into the equation yields $x^2 + y^2 + z^2 = 36$.
\\
\textbf{Final Answer:} This is a sphere centered at the origin with a radius of $6$.
\vspace{0.5cm}

\textbf{Problem 13:} Identify the surface whose equation is given by $z = 5r$.
\\
\textbf{Solution:} Substitute $r = \sqrt{x^2 + y^2}$. The equation becomes $z = 5\sqrt{x^2 + y^2}$, which squares to $z^2 = 25(x^2+y^2)$ for $z \ge 0$.
\\
\textbf{Final Answer:} This is the top half of a cone opening upwards along the $z$-axis.
\vspace{0.5cm}

\textbf{Problem 14:} Identify the surface whose equation is given by $r = 8\sin(\theta)$.
\\
\textbf{Solution:} Multiply both sides by $r$: $r^2 = 8r\sin(\theta)$. Substitute $r^2 = x^2+y^2$ and $y=r\sin(\theta)$:
\begin{align*}
x^2 + y^2 &= 8y \\
x^2 + y^2 - 8y &= 0 \\
x^2 + (y-4)^2 &= 16
\end{align*}
\textbf{Final Answer:} This is a cylinder of radius $4$ centered along the vertical line $(0, 4, z)$.
\vspace{0.5cm}

\textbf{Problem 15:} Express the volume of the solid $E$ bounded by the paraboloid $z = 16 - x^2 - y^2$ and the $xy$-plane ($z = 0$) as an iterated integral in cylindrical coordinates. Do not evaluate.
\\
\textbf{Solution:} Boundaries: Top surface $z = 16 - r^2$, Bottom surface $z = 0$. Intersection: $16 - r^2 = 0 \Rightarrow r = 4$. The region covers a full circle, so $0 \le \theta \le 2\pi$. Integrand is $1$, volume differential is $r \,dz\,dr\,d\theta$.
\\
\textbf{Final Answer:} $\int_0^{2\pi} \int_0^4 \int_0^{16-r^2} r \, dz \, dr \, d\theta$
\vspace{0.5cm}

\textbf{Problem 16:} Express the volume of the solid $E$ bounded below by the cone $z = \sqrt{x^2+y^2}$ and above by the plane $z = 6$ as an iterated integral in cylindrical coordinates. Do not evaluate.
\\
\textbf{Solution:} Boundaries: Bottom surface $z = r$. Top surface $z = 6$. Intersection: $r = 6$. Full circle implies $0 \le \theta \le 2\pi$. Integrand is $1$, differential is $r \,dz\,dr\,d\theta$.
\\
\textbf{Final Answer:} $\int_0^{2\pi} \int_0^6 \int_r^6 r \, dz \, dr \, d\theta$
\vspace{0.5cm}

\textbf{Problem 17:} Evaluate the iterated integral: $\int_0^{2\pi} \int_0^2 \int_0^{r} z r \, dz \, dr \, d\theta$.
\\
\textbf{Solution:} Inner integral with respect to $z$:
\begin{align*}
\int_0^r z r \, dz = r \left[ \frac{z^2}{2} \right]_0^r = r \left( \frac{r^2}{2} \right) = \frac{r^3}{2}
\end{align*}
Middle integral with respect to $r$:
\begin{align*}
\int_0^2 \frac{r^3}{2} \, dr = \left[ \frac{r^4}{8} \right]_0^2 = \frac{16}{8} = 2
\end{align*}
Outer integral with respect to $\theta$:
\begin{align*}
\int_0^{2\pi} 2 \, d\theta = 2(2\pi) = 4\pi
\end{align*}
\textbf{Final Answer:} $4\pi$
\vspace{0.5cm}

\textbf{Problem 18:} Evaluate the iterated integral: $\int_0^{\pi/2} \int_0^3 \int_0^{9-r^2} r \, dz \, dr \, d\theta$.
\\
\textbf{Solution:} Inner integral with respect to $z$:
\begin{align*}
\int_0^{9-r^2} r \, dz = r[z]_0^{9-r^2} = r(9-r^2) = 9r - r^3
\end{align*}
Middle integral with respect to $r$:
\begin{align*}
\int_0^3 (9r - r^3) \, dr = \left[ \frac{9r^2}{2} - \frac{r^4}{4} \right]_0^3 = \frac{81}{2} - \frac{81}{4} = \frac{81}{4}
\end{align*}
Outer integral with respect to $\theta$:
\begin{align*}
\int_0^{\pi/2} \frac{81}{4} \, d\theta = \frac{81}{4} \left(\frac{\pi}{2}\right) = \frac{81\pi}{8}
\end{align*}
\textbf{Final Answer:} $\frac{81\pi}{8}$
\vspace{0.5cm}

\textbf{Problem 19:} Evaluate $\iiint_E (x^2+y^2) \, dV$, where $E$ is the region inside the cylinder $x^2+y^2=4$, and between $z=0$ and $z=3$.
\\
\textbf{Solution:} Convert to cylindrical. Integrand: $x^2+y^2 = r^2$. Bounds: $0 \le z \le 3$, $0 \le r \le 2$, $0 \le \theta \le 2\pi$. Integral is $\int_0^{2\pi} \int_0^2 \int_0^3 (r^2) r \, dz \, dr \, d\theta$.
Inner integral:
\begin{align*}
\int_0^3 r^3 \, dz = r^3[z]_0^3 = 3r^3
\end{align*}
Middle integral:
\begin{align*}
\int_0^2 3r^3 \, dr = \left[ \frac{3r^4}{4} \right]_0^2 = \frac{3(16)}{4} = 12
\end{align*}
Outer integral:
\begin{align*}
\int_0^{2\pi} 12 \, d\theta = 24\pi
\end{align*}
\textbf{Final Answer:} $24\pi$
\vspace{0.5cm}

\textbf{Problem 20:} Evaluate $\iiint_E y \, dV$, where $E$ is the solid in the first octant enclosed by the cylinder $x^2+y^2=9$ and the plane $z=4$.
\\
\textbf{Solution:} First octant means $0 \le \theta \le \pi/2$. Cylinder $r=3$. Bounds: $0 \le z \le 4$, $0 \le r \le 3$. Integrand: $y = r\sin(\theta)$. Differential: $r \, dz \, dr \, d\theta$.
\begin{align*}
\int_0^{\pi/2} \int_0^3 \int_0^4 (r\sin\theta) r \, dz \, dr \, d\theta
\end{align*}
Inner integral:
\begin{align*}
\int_0^4 r^2\sin\theta \, dz = 4r^2\sin\theta
\end{align*}
Middle integral:
\begin{align*}
\int_0^3 4r^2\sin\theta \, dr = \left[ \frac{4r^3}{3} \right]_0^3 \sin\theta = 36\sin\theta
\end{align*}
Outer integral:
\begin{align*}
\int_0^{\pi/2} 36\sin\theta \, d\theta = 36[-\cos\theta]_0^{\pi/2} = 36(0 - (-1)) = 36
\end{align*}
\textbf{Final Answer:} $36$
\vspace{0.5cm}

\textbf{Problem 21:} Find the volume of the solid enclosed by the cone $z = \sqrt{x^2+y^2}$ and the sphere $x^2+y^2+z^2=32$.
\\
\textbf{Solution:} Convert bounds. Cone: $z=r$. Sphere: $z = \sqrt{32-r^2}$. Intersection: $r^2 = 32-r^2 \Rightarrow 2r^2=32 \Rightarrow r=4$.
Integral setup: $V = \int_0^{2\pi} \int_0^4 \int_r^{\sqrt{32-r^2}} r \, dz \, dr \, d\theta$.
Inner integral: $\int_r^{\sqrt{32-r^2}} r \, dz = r(\sqrt{32-r^2} - r) = r\sqrt{32-r^2} - r^2$.
Middle integral split into two parts:
Part 1: $\int_0^4 r\sqrt{32-r^2} \, dr$. Let $u=32-r^2, du=-2rdr$.
\begin{align*}
-\frac{1}{2} \int_{32}^{16} u^{1/2} du &= \frac{1}{2} \left[ \frac{2}{3}u^{3/2} \right]_{16}^{32} = \frac{1}{3}(32\sqrt{32} - 16\sqrt{16}) \\
&= \frac{1}{3}(32(4\sqrt{2}) - 16(4)) = \frac{128\sqrt{2} - 64}{3}
\end{align*}
Part 2: $\int_0^4 r^2 dr = \left[\frac{r^3}{3}\right]_0^4 = \frac{64}{3}$. Subtract Part 2 from Part 1:
\begin{align*}
\frac{128\sqrt{2} - 64}{3} - \frac{64}{3} = \frac{128\sqrt{2} - 128}{3} = \frac{128}{3}(\sqrt{2}-1)
\end{align*}
Outer integral: $\int_0^{2\pi} \frac{128}{3}(\sqrt{2}-1) \, d\theta = \frac{256\pi}{3}(\sqrt{2}-1)$.
\\
\textbf{Final Answer:} $\frac{256\pi}{3}(\sqrt{2}-1)$
\vspace{0.5cm}

\textbf{Problem 22:} Find the volume of the solid that lies within the cylinder $x^2+y^2=9$ and the sphere $x^2+y^2+z^2=25$.
\\
\textbf{Solution:} Cylinder bounds: $0 \le r \le 3$. Sphere bounds: $z = \pm\sqrt{25-r^2}$.
Integral: $\int_0^{2\pi} \int_0^3 \int_{-\sqrt{25-r^2}}^{\sqrt{25-r^2}} r \, dz \, dr \, d\theta$.
Inner integral: $r [z] = 2r\sqrt{25-r^2}$.
Middle integral: $\int_0^3 2r\sqrt{25-r^2} \, dr$. Use $u = 25-r^2$, $du = -2rdr$. Limits: $25$ to $16$.
\begin{align*}
-\int_{25}^{16} u^{1/2} du = \int_{16}^{25} u^{1/2} du = \left[ \frac{2}{3}u^{3/2} \right]_{16}^{25} = \frac{2}{3}(125 - 64) = \frac{2}{3}(61) = \frac{122}{3}
\end{align*}
Outer integral: $\int_0^{2\pi} \frac{122}{3} d\theta = \frac{244\pi}{3}$.
\\
\textbf{Final Answer:} $\frac{244\pi}{3}$
\vspace{0.5cm}

\textbf{Problem 23:} Evaluate $\iiint_E z \, dV$, where $E$ is the region in the first octant bounded by the cylinder $x^2+y^2=16$ and the planes $z=0$, $z=y$, and $x=0$.
\\
\textbf{Solution:} First octant bounds: $0 \le \theta \le \pi/2$. Cylinder $r=4$. Top plane $z=y \Rightarrow z=r\sin\theta$.
$\int_0^{\pi/2} \int_0^4 \int_0^{r\sin\theta} z r \, dz \, dr \, d\theta$.
Inner:
\begin{align*}
\int_0^{r\sin\theta} zr \, dz = r\left[\frac{z^2}{2}\right]_0^{r\sin\theta} = \frac{r^3}{2}\sin^2\theta
\end{align*}
Middle:
\begin{align*}
\int_0^4 \frac{r^3}{2}\sin^2\theta \, dr = \left[ \frac{r^4}{8} \right]_0^4 \sin^2\theta = 32\sin^2\theta
\end{align*}
Outer:
\begin{align*}
\int_0^{\pi/2} 32\sin^2\theta \, d\theta = 32 \int_0^{\pi/2} \frac{1-\cos(2\theta)}{2} d\theta = 16 \left[ \theta - \frac{\sin(2\theta)}{2} \right]_0^{\pi/2} = 16\left(\frac{\pi}{2} - 0\right) = 8\pi
\end{align*}
\textbf{Final Answer:} $8\pi$
\vspace{0.5cm}

\textbf{Problem 24:} Evaluate $\iiint_E x^2 \, dV$, where $E$ is the solid inside the cylinder $x^2+y^2=4$ bounded by $z=0$ and $z=5$.
\\
\textbf{Solution:} $x^2 = r^2\cos^2\theta$. Full circle $0 \le \theta \le 2\pi$. $0 \le r \le 2$. Integral: $\int_0^{2\pi} \int_0^2 \int_0^5 (r^2\cos^2\theta) r \, dz \, dr \, d\theta$.
Inner: $\int_0^5 r^3\cos^2\theta \, dz = 5r^3\cos^2\theta$.
Middle: $\int_0^2 5r^3\cos^2\theta \, dr = \left[ \frac{5r^4}{4} \right]_0^2 \cos^2\theta = 20\cos^2\theta$.
Outer: 
\begin{align*}
\int_0^{2\pi} 20\cos^2\theta \, d\theta = 20 \int_0^{2\pi} \frac{1+\cos(2\theta)}{2} d\theta = 10\left[\theta + \frac{\sin(2\theta)}{2}\right]_0^{2\pi} = 20\pi
\end{align*}
\textbf{Final Answer:} $20\pi$
\vspace{0.5cm}

\textbf{Problem 25:} Evaluate $\iiint_E (y^3 + z) \, dV$, where $E$ is bounded by $x^2+y^2=1$, $z=0$, and $z=2$.
\\
\textbf{Solution:} By independence and odd symmetry, the integral of $y^3$ over a full symmetric domain $[0, 2\pi]$ is $0$ ($\int_0^{2\pi} \sin^3\theta d\theta = 0$). So, we only evaluate $\iiint_E z \, dV$.
Integral: $\int_0^{2\pi} \int_0^1 \int_0^2 z r \, dz \, dr \, d\theta$.
Inner: $\int_0^2 zr \, dz = r[z^2/2]_0^2 = 2r$.
Middle: $\int_0^1 2r \, dr = [r^2]_0^1 = 1$.
Outer: $\int_0^{2\pi} 1 \, d\theta = 2\pi$.
\\
\textbf{Final Answer:} $2\pi$
\vspace{0.5cm}

\textbf{Problem 26:} Evaluate $\iiint_E (x - y + 4) \, dV$, where $E$ lies inside the cylinder $x^2+y^2=9$, $z=0$, and $z=1$.
\\
\textbf{Solution:} Apply the Fubini symmetry trick. Since $x = r\cos\theta$ and $y = r\sin\theta$, integrating these to the first power over a full $2\pi$ period yields zero. We are left only with $\iiint_E 4 \, dV$.
This equals $4 \times (\text{Volume of } E)$.
The volume of a cylinder is $\pi r^2 h = \pi (3)^2 (1) = 9\pi$.
$4 \times 9\pi = 36\pi$.
\\
\textbf{Final Answer:} $36\pi$
\vspace{0.5cm}

\textbf{Problem 27:} Find the mass of the solid cylinder $E$ defined by $0 \le r \le 2$ and $0 \le z \le 3$, whose density is given by $\rho(x,y,z) = z\sqrt{x^2+y^2}$.
\\
\textbf{Solution:} Density $\rho = z \cdot r$. Mass $m = \iiint_E \rho \, dV$.
Integral: $m = \int_0^{2\pi} \int_0^2 \int_0^3 (zr) r \, dz \, dr \, d\theta$.
Inner: $\int_0^3 z r^2 \, dz = r^2\left[\frac{z^2}{2}\right]_0^3 = \frac{9}{2}r^2$.
Middle: $\int_0^2 \frac{9}{2}r^2 \, dr = \frac{9}{2}\left[\frac{r^3}{3}\right]_0^2 = \frac{3}{2}(8) = 12$.
Outer: $\int_0^{2\pi} 12 \, d\theta = 24\pi$.
\\
\textbf{Final Answer:} $24\pi$
\vspace{0.5cm}

\textbf{Problem 28:} Find the mass of the solid $E$ bounded by the paraboloid $z = x^2+y^2$ and the plane $z=4$, given that the density of the solid is a constant $\rho(x,y,z) = 2$.
\\
\textbf{Solution:} Paraboloid $z=r^2$. Intersection $r^2=4 \Rightarrow r=2$.
Mass $m = \int_0^{2\pi} \int_0^2 \int_{r^2}^4 2 r \, dz \, dr \, d\theta$.
Inner: $\int_{r^2}^4 2r \, dz = 2r(4 - r^2) = 8r - 2r^3$.
Middle: $\int_0^2 (8r - 2r^3) \, dr = \left[ 4r^2 - \frac{r^4}{2} \right]_0^2 = 16 - \frac{16}{2} = 8$.
Outer: $\int_0^{2\pi} 8 \, d\theta = 16\pi$.
\\
\textbf{Final Answer:} $16\pi$
\vspace{0.5cm}

\textbf{Problem 29:} \textit{Find the Flaw.} A student is tasked with evaluating the volume inside the cylinder $x^2+y^2=9$ and between the planes $z=-2$ and $z=2$. The student sets up the following integral: $\int_0^{2\pi}\int_0^3\int_{-2}^2 1 \,dz\,dr\,d\theta$. Identify the error, explain why it is wrong, and provide the correct calculation.
\\
\textbf{Solution:} \textbf{Flaw:} The student omitted the Jacobian $r$ in the integrand. The volume differential must be $dV = r \, dz \, dr \, d\theta$, so the integrand should be $r$, not $1$.
\textbf{Correct Calculation:} $V = \int_0^{2\pi}\int_0^3\int_{-2}^2 r \,dz\,dr\,d\theta$. Inner: $\int_{-2}^2 r \, dz = 4r$. Middle: $\int_0^3 4r \, dr = 18$. Outer: $\int_0^{2\pi} 18 \, d\theta = 36\pi$.
\\
\textbf{Final Answer:} The flaw is the missing Jacobian $r$. Correct volume is $36\pi$.
\vspace{0.5cm}

\textbf{Problem 30:} \textit{Find the Flaw.} A student is setting up an integral to find the volume of the solid enclosed between the cone $z = 3r$ and the $xy$-plane ($z=0$) inside the cylinder $r=2$. The student writes: $V = \int_0^{2\pi}\int_0^2\int_{3r}^0 r \,dz\,dr\,d\theta$. Identify the flaw, explain why it is wrong, and provide the correct evaluation.
\\
\textbf{Solution:} \textbf{Flaw:} The student reversed the $z$ bounds, placing $3r$ on the bottom and $0$ on top. This forces the evaluation of a negative volume. Since the solid lies between the cone and the $xy$-plane, the region is bounded below by $z=0$ and above by $z=3r$. 
\textbf{Correct Calculation:} $V = \int_0^{2\pi}\int_0^2\int_0^{3r} r \,dz\,dr\,d\theta$. Inner: $\int_0^{3r} r \, dz = r(3r) = 3r^2$. Middle: $\int_0^2 3r^2 \, dr = [r^3]_0^2 = 8$. Outer: $\int_0^{2\pi} 8 \, d\theta = 16\pi$.
\\
\textbf{Final Answer:} The flaw is the reversed $z$ bounds ($\int_{3r}^0$ instead of $\int_0^{3r}$). Correct volume is $16\pi$.
\vspace{1cm}

\pagebreak

\section*{Concept Checklist}
Below is the mapping of concepts and problem types tested in this homework assignment to their corresponding question numbers:

\begin{itemize}
    \item \textbf{Concept 1: Point Conversion (Cylindrical to Rectangular)} \\ \textit{Questions: 1, 2, 3, 4}
    \item \textbf{Concept 2: Point Conversion (Rectangular to Cylindrical)} \\ \textit{Questions: 5, 6, 7, 8}
    \item \textbf{Concept 3: Equation Conversion (Rectangular to Cylindrical)} \\ \textit{Questions: 9, 10, 11}
    \item \textbf{Concept 4: Surface Identification} \\ \textit{Questions: 12, 13, 14}
    \item \textbf{Concept 5: Boundary Setup Problems (No Evaluation)} \\ \textit{Questions: 15, 16}
    \item \textbf{Concept 6: Full Iterated Evaluation (Basic to Moderate)} \\ \textit{Questions: 17, 18, 19, 20}
    \item \textbf{Concept 7: U-Substitution and Fractional Powers (Spheres/Cones)} \\ \textit{Questions: 21, 22}
    \item \textbf{Concept 8: Double Angle Identity for Trig Powers} \\ \textit{Questions: 23, 24}
    \item \textbf{Concept 9: Fubini Simplification Trick and Symmetry Shortcuts} \\ \textit{Questions: 25, 26}
    \item \textbf{Concept 10: Applications (Mass and Variable Density)} \\ \textit{Questions: 27, 28}
    \item \textbf{Concept 11: Diagnostic Testing ("Find the Flaw" and Pitfalls)} \\ \textit{Questions: 29, 30}
\end{itemize}

\end{document}